\subsection{Dạng 1. Tính số đo góc (bằng cách vẽ thêm hình)}

\begin{bt}
	Cho hình vẽ, biết $MN\parallel KJ$,
	$\widehat{NML}=46^\circ$ và $\widehat{JKL}=127^\circ$.
	Hãy tính số đo góc $\widehat{MLK}$.

	\begin{center}
		\begin{tikzpicture}[
			scale=1.1,
			line cap=round,
			line join=round,
			font=\small,
			every node/.style={inner sep=1.3pt}
		]
			\coordinate (M) at (0,1.5);
			\coordinate (N) at (3.05,1.5);
			\coordinate (L) at ($(M)+(-46:2.05)$);
			\coordinate (K) at ($(L)+(233:1.25)$);
			\coordinate (J) at ($(K)+(-2.35,0)$);

			% Hai đường thẳng song song và đường gấp MLK
			\draw[mauChuDao,line width=0.9pt] (M)--(N);
			\draw[mauChuDao,line width=0.9pt] (J)--(K);
			\draw[mauChuDao,line width=0.9pt] (M)--(L)--(K);

			% Hai góc đã biết
			\pic[draw=mauChuDao,"$46^\circ$",angle radius=5mm,
				angle eccentricity=1.55] {angle=L--M--N};
			\pic[draw=mauChuDao,"$127^\circ$",angle radius=5mm,
				angle eccentricity=1.45] {angle=L--K--J};

			% Các điểm và nhãn
			\foreach \Pnt in {M,N,L,K,J}{
				\fill[red!75!black] (\Pnt) circle (1.45pt);
			}
			\node[above left=2pt] at (M) {$M$};
			\node[above right=2pt] at (N) {$N$};
			\node[right=3pt] at (L) {$L$};
			\node[below right=1pt] at (K) {$K$};
			\node[below left=1pt] at (J) {$J$};
		\end{tikzpicture}
	\end{center}
\end{bt}

\begin{bt}
	Cho hình vẽ, biết $AB\parallel ED$,
	$\widehat{BAC}=118^\circ$ và $\widehat{CDE}=50^\circ$.
	Hãy tính số đo góc $\widehat{ACD}$.

	\begin{center}
		\begin{tikzpicture}[
			scale=1.1,
			line cap=round,
			line join=round,
			font=\small,
			every node/.style={inner sep=1.3pt}
		]
			\coordinate (A) at (0,1.55);
			\coordinate (B) at (2.8,1.55);
			\coordinate (C) at ($(A)+(242:1.45)$);
			\coordinate (D) at ($(C)+(310:1.55)$);
			\coordinate (E) at ($(D)+(-2.85,0)$);

			% Hai đường thẳng song song và đường gấp ACD
			\draw[mauChuDao,line width=0.9pt] (A)--(B);
			\draw[mauChuDao,line width=0.9pt] (E)--(D);
			\draw[mauChuDao,line width=0.9pt] (A)--(C)--(D);

			% Hai góc đã biết
			\pic[draw=mauChuDao,"$118^\circ$",angle radius=4.8mm,
				angle eccentricity=1.45] {angle=C--A--B};
			\pic[draw=mauChuDao,"$50^\circ$",angle radius=4.8mm,
				angle eccentricity=1.5] {angle=C--D--E};

			% Các điểm và nhãn
			\foreach \Pnt in {A,B,C,D,E}{
				\fill[red!75!black] (\Pnt) circle (1.45pt);
			}
			\node[above left=2pt] at (A) {$A$};
			\node[above right=2pt] at (B) {$B$};
			\node[left=3pt] at (C) {$C$};
			\node[below right=1pt] at (D) {$D$};
			\node[below left=1pt] at (E) {$E$};
		\end{tikzpicture}
	\end{center}
\end{bt}

\begin{bt}
	Cho hình vẽ, biết ba điểm $C,B,E$ thẳng hàng, $AB\parallel FG$,
	$\widehat{ABC}=49^\circ$ và $\widehat{EFG}=120^\circ$.
	Hãy tính số đo góc $\widehat{CEF}$.

	\begin{center}
		\begin{tikzpicture}[
			scale=1.1,
			line cap=round,
			line join=round,
			font=\small,
			every node/.style={inner sep=1.3pt}
		]
			\coordinate (B) at (0,0.9);
			\coordinate (A) at (2.8,0.9);
			\coordinate (C) at ($(B)+(49:1.5)$);
			\coordinate (E) at ($(B)+(229:0.85)$);
			\coordinate (F) at ($(E)+(300:1.3)$);
			\coordinate (G) at ($(F)+(2.8,0)$);

			% Hai đường thẳng song song và đường gấp CEF
			\draw[mauChuDao,line width=0.9pt] (B)--(A);
			\draw[mauChuDao,line width=0.9pt] (F)--(G);
			\draw[mauChuDao,line width=0.9pt] (C)--(B)--(E)--(F);

			% Hai góc đã biết
			\pic[draw=mauChuDao,"$49^\circ$",angle radius=4.7mm,
				angle eccentricity=1.7] {angle=A--B--C};
			\pic[draw=mauChuDao,"$120^\circ$",angle radius=3.8mm,
				angle eccentricity=1.7] {angle=G--F--E};

			% Các điểm và nhãn
			\foreach \Pnt in {A,B,C,E,F,G}{
				\fill[red!75!black] (\Pnt) circle (1.45pt);
			}
			\node[above right=2pt] at (A) {$A$};
			\node[above left=1pt] at (B) {$B$};
			\node[above right=1pt] at (C) {$C$};
			\node[above left=1pt] at (E) {$E$};
			\node[below left=1pt] at (F) {$F$};
			\node[below right=1pt] at (G) {$G$};
		\end{tikzpicture}
	\end{center}
\end{bt}

\subsection{Dạng 2. Chứng minh hai đường thẳng song song}

\begin{bt}
	Theo hình vẽ, cho
	$\widehat{xAB}+\widehat{ABC}+\widehat{BCy}=360^\circ$.
	Hãy chứng minh $Ax\parallel Cy$.

	\begin{center}
		\begin{tikzpicture}[
			scale=1.1,
			line cap=round,
			line join=round,
			font=\small,
			every node/.style={inner sep=1.3pt}
		]
			\coordinate (A) at (1.3,1.35);
			\coordinate (B) at (2.75,0);
			\coordinate (C) at (1.3,-1.25);
			\coordinate (X) at (-2.7,1.35);
			\coordinate (Y) at (-3.35,-1.25);

			% Hai tia Ax, Cy và đường gấp ABC
			\draw[mauChuDao,line width=0.9pt] (X)--(A)--(B)--(C)--(Y);

			% Các đỉnh của đường gấp
			\foreach \Pnt in {A,B,C}{
				\fill[red!75!black] (\Pnt) circle (1.45pt);
			}

			% Nhãn các điểm và tên tia
			\node[above right=2pt] at (A) {$A$};
			\node[right=4pt] at (B) {$B$};
			\node[below right=2pt] at (C) {$C$};
			\node[above right=1pt] at (X) {$x$};
			\node[above right=1pt] at (Y) {$y$};
		\end{tikzpicture}
	\end{center}
\end{bt}

\begin{bt}
	Cho hình vẽ bên. Biết $My\parallel Pt$,
	$\widehat{xMy}=60^\circ$ và $\widehat{MNz}=60^\circ$.
	\begin{enumerate}[label=\alph*)]
		\item Chứng minh $My\parallel Nz$;
		\item Chứng minh $Nz\parallel Pt$;
		\item Biết $\widehat{xNP}=85^\circ$, tính số đo góc $\widehat{NPt}$.
	\end{enumerate}

	\begin{center}
		\begin{tikzpicture}[
			scale=1.1,
			line cap=round,
			line join=round,
			font=\small,
			every node/.style={inner sep=1.3pt}
		]
			\coordinate (N) at (0,0);
			\coordinate (M) at ($(N)+(60:1.15)$);
			\coordinate (X) at ($(M)+(60:1.2)$);
			\coordinate (Y) at ($(M)+(2.45,0)$);
			\coordinate (Z) at ($(N)+(2.55,0)$);
			\coordinate (P) at ($(N)+(-25:1.65)$);
			\coordinate (T) at ($(P)+(-2.35,0)$);

			% Các tia và đoạn thẳng trong hình
			\draw[mauChuDao,line width=0.9pt] (N)--(M)--(X);
			\draw[mauChuDao,line width=0.9pt] (M)--(Y);
			\draw[mauChuDao,line width=0.9pt] (N)--(Z);
			\draw[mauChuDao,line width=0.9pt] (N)--(P)--(T);

			% Hai góc 60 độ
			\pic[draw=mauChuDao,"$60^\circ$",angle radius=3.5mm,
				angle eccentricity=1.65] {angle=Y--M--X};
			\pic[draw=mauChuDao,"$60^\circ$",angle radius=3.5mm,
				angle eccentricity=1.65] {angle=Z--N--M};

			% Các điểm
			\foreach \Pnt in {M,N,P}{
				\fill[red!75!black] (\Pnt) circle (1.45pt);
			}

			% Nhãn các điểm và tên tia
			\node[left=3pt] at (M) {$M$};
			\node[left=3pt] at (N) {$N$};
			\node[below right=1pt] at (P) {$P$};
			\node[above left=1pt] at (X) {$x$};
			\node[below right=1pt] at (Y) {$y$};
			\node[below right=1pt] at (Z) {$z$};
			\node[below left=1pt] at (T) {$t$};
		\end{tikzpicture}
	\end{center}
\end{bt}

\begin{bt}
	Cho hình vẽ sau, biết $HD\parallel GE$ và $GE\parallel IF$.
	\begin{enumerate}[label=\alph*)]
		\item Chứng minh $HD\parallel IF$.
		\item Cho biết $\widehat{HAB}=65^\circ$. Tìm số đo $x$ của góc
		$\widehat{ABE}$ và số đo $y$ của góc $\widehat{BCI}$.
	\end{enumerate}

	\begin{center}
		\begin{tikzpicture}[
			scale=1.05,
			line cap=round,
			line join=round,
			font=\small,
			every node/.style={inner sep=1.3pt}
		]
			\coordinate (A) at (0,1.9);
			\coordinate (B) at ($(A)+(245:1.15)$);
			\coordinate (C) at ($(B)+(245:1.15)$);
			\coordinate (T) at ($(A)+(65:0.7)$);
			\coordinate (K) at ($(C)+(245:0.65)$);

			\coordinate (H) at (-2.25,1.9);
			\coordinate (D) at (1.75,1.9);
			\coordinate (G) at (-2.25,{1.9-1.15*sin(65)});
			\coordinate (E) at (1.75,{1.9-1.15*sin(65)});
			\coordinate (I) at (-2.25,{1.9-2.3*sin(65)});
			\coordinate (F) at (1.75,{1.9-2.3*sin(65)});

			% Ba đường thẳng song song và cát tuyến
			\draw[mauChuDao,line width=0.9pt] (H)--(D);
			\draw[mauChuDao,line width=0.9pt] (G)--(E);
			\draw[mauChuDao,line width=0.9pt] (I)--(F);
			\draw[mauChuDao,line width=0.9pt] (K)--(T);

			% Góc HAB và hai góc cần tìm x, y
			\pic[draw=mauChuDao,"$65^\circ$",angle radius=3.5mm,
				angle eccentricity=1.5] {angle=H--A--B};
			\pic[draw=mauChuDao,"$x$",angle radius=3.2mm,
				angle eccentricity=1.55] {angle=E--B--A};
			\pic[draw=mauChuDao,"$y$",angle radius=4.5mm,
				angle eccentricity=1.5] {angle=B--C--I};

			% Các giao điểm
			\foreach \Pnt in {A,B,C}{
				\fill[red!75!black] (\Pnt) circle (1.45pt);
			}

			% Nhãn các điểm
			\node[left=3pt] at (H) {$H$};
			\node[right=3pt] at (D) {$D$};
			\node[left=3pt] at (G) {$G$};
			\node[right=3pt] at (E) {$E$};
			\node[left=3pt] at (I) {$I$};
			\node[right=3pt] at (F) {$F$};
			\node[above left=1pt] at (A) {$A$};
			\node[above left=8pt] at (B) {$B$};
			\node[below right=1pt] at (C) {$C$};
			\node[above right=1pt] at (T) {$T$};
			\node[below left=1pt] at (K) {$K$};
		\end{tikzpicture}
	\end{center}
\end{bt}
\begin{bt}
	Cho hình vẽ, biết $\widehat{A_2}=\widehat{M_1}$ và
	$\widehat{A_1}=\widehat{N_2}$.
	\begin{enumerate}[label=\alph*)]
		\item Cho biết $\widehat{A_2}$ và $\widehat{M_1}$ là hai góc ở vị trí
		như thế nào.
		\item Hai góc $\widehat{A_1}$ và $\widehat{N_2}$ là hai góc gì?
		\item Chứng minh $MD\parallel AH$.
		\item Chứng minh $NE\parallel MD$.
	\end{enumerate}

	\begin{center}
		\begin{tikzpicture}[
			scale=1.05,
			line cap=round,
			line join=round,
			font=\small,
			every node/.style={inner sep=1.2pt}
		]
			\coordinate (A) at (0,3.0);
			\coordinate (H) at (0,0);
			\coordinate (B) at (-2.9,0);
			\coordinate (C) at (1.55,0);
			\coordinate (N) at (2.2475,-1.35);
			\coordinate (E) at (2.2475,0);
			\coordinate (M) at ($(B)!0.52!(A)$);
			\coordinate (D) at ($(M |- H)$);

			% Đường đáy và các đường trong hình
			\draw[mauChuDao,line width=0.9pt] (B)--(E);
			\draw[mauChuDao,line width=0.9pt] (B)--(A)--(C)--(N);
			\draw[mauChuDao,line width=0.9pt] (A)--(H);
			\draw[mauChuDao,line width=0.9pt] (M)--(D);
			\draw[mauChuDao,line width=0.9pt] (N)--(E);

			% Ký hiệu các góc A_1, A_2, M_1 và N_2
			\node[font=\scriptsize] at ($(A)+(-76:0.30)$) {$1$};
			\node[font=\scriptsize] at ($(A)+(-112:0.30)$) {$2$};
			\node[font=\scriptsize] at ($(M)+(-112:0.28)$) {$1$};
			\node[font=\scriptsize] at ($(N)+(103:0.30)$) {$2$};

			\foreach \Pnt in {A,B,C,D,E,H,M,N}{
				\fill[red!75!black] (\Pnt) circle (1.35pt);
			}
			\node[above=3pt] at (A) {$A$};
			\node[below left=1pt] at (B) {$B$};
			\node[above left=1pt] at (C) {$C$};
			\node[below=3pt] at (D) {$D$};
			\node[above right=1pt] at (E) {$E$};
			\node[below=3pt] at (H) {$H$};
			\node[left=3pt] at (M) {$M$};
			\node[below right=1pt] at (N) {$N$};
		\end{tikzpicture}
	\end{center}
\end{bt}

\subsection{Dạng 3. Chứng minh song song do cùng vuông góc}

\begin{bt}
	Cho góc nhọn $\widehat{AOB}$. Từ $M$ trên tia $OA$ vẽ
	$MN\perp OB$ $(N\in OB)$; từ $N$ vẽ $NP\perp OA$ $(P\in OA)$;
	từ $P$ vẽ $PQ\perp OB$ $(Q\in OB)$; từ $Q$ vẽ
	$QR\perp OA$ $(R\in OA)$.
	\begin{enumerate}[label=\alph*)]
		\item Chứng minh $MN\parallel PQ$ và $NP\parallel QR$.
		\item Trong hình vẽ, xác định tất cả các góc có số đo bằng
		$\widehat{PMN}$ và tất cả các góc có số đo bằng $\widehat{MNP}$.
	\end{enumerate}

	\begin{center}
		\begin{tikzpicture}[
			scale=1.05,
			line cap=round,
			line join=round,
			font=\small,
			every node/.style={inner sep=1.3pt}
		]
			\coordinate (O) at (0,0);
			\coordinate (B) at (5.0,0);
			\coordinate (A) at (4.05,3.16);
			\coordinate (M) at (3.395,2.648);
			\coordinate (N) at (3.395,0);
			\coordinate (P) at (2.108,1.647);
			\coordinate (Q) at (2.108,0);
			\coordinate (R) at (1.309,1.023);

			% Hai tia của góc AOB
			\draw[mauChuDao,line width=0.9pt] (O)--(A);
			\draw[mauChuDao,line width=0.9pt] (O)--(B);

			% Các đoạn vuông góc liên tiếp
			\draw[mauChuDao,line width=0.9pt] (M)--(N)--(P)--(Q)--(R);

			% Dấu vuông góc tại N, P, Q và R
			\pic[draw=mauChuDao,angle radius=2.5mm]
				{right angle=M--N--O};
			\pic[draw=mauChuDao,angle radius=2.5mm]
				{right angle=M--P--N};
			\pic[draw=mauChuDao,angle radius=2.5mm]
				{right angle=P--Q--O};
			\pic[draw=mauChuDao,angle radius=2.5mm]
				{right angle=P--R--Q};

			% Các điểm
			\foreach \Pnt in {O,M,N,P,Q,R}{
				\fill[red!75!black] (\Pnt) circle (1.4pt);
			}

			% Nhãn các điểm
			\node[below left=1pt] at (O) {$O$};
			\node[above left=1pt] at (A) {$A$};
			\node[below right=1pt] at (B) {$B$};
			\node[above left=1pt] at (M) {$M$};
			\node[below right=1pt] at (N) {$N$};
			\node[above left=1pt] at (P) {$P$};
			\node[below=3pt] at (Q) {$Q$};
			\node[above left=1pt] at (R) {$R$};
		\end{tikzpicture}
	\end{center}
\end{bt}

\subsection{Dạng 4. Chứng minh vuông góc từ song song}

\begin{bt}
	Cho hình vẽ. Biết $a\parallel b$, $b\perp c$ và tia $Ad$ là tia
	phân giác của góc $\widehat{cAx}$, trong đó tia $Ax$ nằm trên đường
	thẳng $a$.
	\begin{enumerate}[label=\alph*)]
		\item Chứng minh $a\perp c$.
		\item Tính số đo góc $\widehat{cAd}$.
	\end{enumerate}

	\begin{center}
		\begin{tikzpicture}[
			scale=1.05,
			line cap=round,
			line join=round,
			font=\small,
			every node/.style={inner sep=1.3pt}
		]
			\coordinate (A) at (0,1.05);
			\coordinate (B) at (0,-1.0);
			\coordinate (X) at (2.45,1.05);
			\coordinate (D) at ($(A)+(45:1.55)$);
			\coordinate (BR) at (1,-1.0);
			\coordinate (CT) at (0,2.05);

			% Hai đường thẳng song song a, b và đường thẳng c
			\draw[mauChuDao,line width=0.9pt]
				(-1.35,1.05)--(2.6,1.05) node[right] {$a$};
			\draw[mauChuDao,line width=0.9pt]
				(-1.35,-1.0)--(2.6,-1.0) node[right] {$b$};
			\draw[mauChuDao,line width=0.9pt]
				(0,-1.45)--(0,2.35) node[above] {$c$};
			\draw[mauChuDao,line width=0.9pt] (A)--(D)
				node[above right] {$d$};

			% Dấu vuông tại B và hai cung bằng nhau tại A
			\pic[draw=mauChuDao,angle radius=2.7mm]
				{right angle=A--B--BR};
			\pic[draw=mauChuDao,angle radius=4.5mm]
				{angle=D--A--CT};
			\pic[draw=mauChuDao,angle radius=5.2mm]
				{angle=X--A--D};

			\fill[red!75!black] (A) circle (1.45pt);
			\fill[red!75!black] (B) circle (1.45pt);
			\node[below left=1pt] at (A) {$A$};
			\node[below left=1pt] at (B) {$B$};
			\node[above right=1pt] at (X) {$x$};
		\end{tikzpicture}
	\end{center}
\end{bt}

\begin{bt}
	Cho hình vẽ. Đường thẳng $c$ cắt $m$ tại $C$ và cắt $n$ tại $D$.
	Biết $\widehat{C_1}=\widehat{D_1}$ và $AB\perp m$.
	\begin{enumerate}[label=\alph*)]
		\item Chứng minh $m\parallel n$.
		\item Chứng minh $AB\perp n$.
	\end{enumerate}

	\begin{center}
		\begin{tikzpicture}[
			scale=1.05,
			line cap=round,
			line join=round,
			font=\small,
			every node/.style={inner sep=1.3pt}
		]
			\coordinate (A) at (-1.45,1.0);
			\coordinate (B) at (-1.45,-1.0);
			\coordinate (C) at (0.35,1.0);
			\coordinate (D) at (1.05,-1.0);
			\coordinate (AR) at (-0.45,1.0);
			\coordinate (CL) at (-0.65,1.0);
			\coordinate (DR) at (2.05,-1.0);

			% Hai đường m, n; đường AB và cát tuyến c
			\draw[mauChuDao,line width=0.9pt]
				(-2.25,1.0)--(2.6,1.0) node[right] {$m$};
			\draw[mauChuDao,line width=0.9pt]
				(-2.25,-1.0)--(2.6,-1.0) node[right] {$n$};
			\draw[mauChuDao,line width=0.9pt] (-1.45,-1.4)--(-1.45,1.4);
			\draw[mauChuDao,line width=0.9pt]
				($(C)!-0.3!(D)$)--($(C)!1.3!(D)$) node[below right] {$c$};

			% Dấu vuông tại A
			\pic[draw=mauChuDao,angle radius=2.7mm]
				{right angle=B--A--AR};

			% Hai góc so le trong bằng nhau
			\pic[draw=mauChuDao,"$1$",angle radius=4.5mm,
				angle eccentricity=1.45]
				{angle=CL--C--D};
			\pic[draw=mauChuDao,"$1$",angle radius=4.5mm,
				angle eccentricity=1.45]
				{angle=DR--D--C};

			% Các giao điểm và nhãn
			\foreach \Pnt in {A,B,C,D}{
				\fill[red!75!black] (\Pnt) circle (1.45pt);
			}
			\node[above left=1pt] at (A) {$A$};
			\node[below left=1pt] at (B) {$B$};
			\node[above right=1pt] at (C) {$C$};
			\node[below left=1pt] at (D) {$D$};
		\end{tikzpicture}
	\end{center}
\end{bt}

\subsection{Dạng 5. Chứng minh thẳng hàng}

\begin{bt}
	Cho hình vẽ bên, biết $Ax\parallel Dy$,
	$\widehat{BCD}=135^\circ$, $AE\perp Ax$ và $CD\perp Dy$.
	\begin{enumerate}[label=\arabic*)]
		\item Đường thẳng $AE$ có vuông góc với $Dy$ không? Vì sao?
		\item Chứng minh $DC\parallel AE$.
		\item Vẽ tia phân giác $Dt$ của góc $\widehat{CDy}$; tia $Dt$ cắt
		đường thẳng $AE$ tại $F$. Tính $\widehat{EDF}$ và $\widehat{DFB}$.
		\item Trên nửa mặt phẳng có bờ là đường thẳng $AB$ không chứa điểm
		$C$, vẽ tia $Bz$ sao cho $\widehat{ABz}=45^\circ$. Tia $Bz$ cắt
		tia $Ax$ tại $M$. Chứng minh ba điểm $B,C,M$ thẳng hàng.
	\end{enumerate}

	\begin{center}
		\begin{tikzpicture}[
			scale=1.0,
			line cap=round,
			line join=round,
			font=\small,
			every node/.style={inner sep=1.3pt}
		]
			\coordinate (D) at (0,0);
			\coordinate (C) at (0,1.5);
			\coordinate (E) at (1.5,0);
			\coordinate (B) at (1.5,3.0);
			\coordinate (A) at (1.5,3.75);
			\coordinate (X) at (4.25,3.75);
			\coordinate (Y) at (4.25,0);
			\coordinate (ER) at (2.5,0);
			\coordinate (AD) at (1.5,2.75);

			% Hai tia song song Ax, Dy và các đoạn AE, CD, CB
			\draw[mauChuDao,line width=0.9pt] (A)--(X);
			\draw[mauChuDao,line width=0.9pt] (D)--(Y);
			\draw[mauChuDao,line width=0.9pt] (A)--(E);
			\draw[mauChuDao,line width=0.9pt] (D)--(C)--(B);

			% Hai dấu vuông góc tại A và D
			\pic[draw=mauChuDao,angle radius=3mm]
				{right angle=AD--A--X};
			\pic[draw=mauChuDao,angle radius=3mm]
				{right angle=C--D--ER};

			% Góc BCD bằng 135 độ
			\pic[draw=mauChuDao,"$135^\circ$",angle radius=2.5mm,
				angle eccentricity=2.5] {angle=D--C--B};

			% Các điểm
			\foreach \Pnt in {A,B,C,D,E}{
				\fill[red!75!black] (\Pnt) circle (1.5pt);
			}

			% Nhãn các điểm và tên tia
			\node[above left=1pt] at (A) {$A$};
			\node[left=3pt] at (B) {$B$};
			\node[left=3pt] at (C) {$C$};
			\node[below left=1pt] at (D) {$D$};
			\node[below right=1pt] at (E) {$E$};
			\node[below left=1pt] at (X) {$x$};
			\node[below left=1pt] at (Y) {$y$};
		\end{tikzpicture}
	\end{center}
\end{bt}



